% 第9回　動線とUXのデザイン（記入例・模範解答）
\session{9}{2}{11月25日（水）}{動線とUXのデザイン}{AIとシミュレーション}{yes}

\goals{
  \item 来店から退店までの顧客体験を、行動、感情、接点に分けて記述できる
  \item AIをペルソナとして動かし、店内行動を予測させる方法を実践できる
  \item 予測から問題点を抽出し、レイアウトの改善案として具体化できる
  \item AIの予測は仮説であり、実際の観察や企業の知見で確かめる必要があることを理解する
}

\work[80mm]{ワーク1}{カスタマージャーニーマップ}
\afillin{ペルソナ／来店目的}{ペルソナA（森 由佳・子連れ）／ミニバンへの買い替えの下見}
\begin{wstabh}{colspec={Q[l,wd=11mm,m] X[c,m] X[c,m] X[c,m] X[c,m] X[c,m] X[c,m] X[c,m]},
    row{2-Z}={ht=19mm},column{1}={font=\sffamily\footnotesize\bfseries}}
  & 来店前 & 入店 & 展示を見る & 商談 & 待ち時間 & 退店 & 退店後 \\
  行動 & \af{サイトで車種と店の雰囲気を見る} & \af{子どもと入店。カードを受け取る} & \af{子どもをキッズスペースへ。展示車に乗る} & \af{カードで相談。見積もり} & \af{見積もり作成を待つ} & \af{資料を受け取り帰る} & \af{家で夫と比べる} \\
  感情\newline\scriptsize ＋／− & \af{＋ 子連れ歓迎とわかり安心} & \af{＋ 声をかけられず気が楽} & \af{− 子どもが見えない位置だと不安} & \af{＋ 横並びで話しやすい} & \af{− 子どもが飽きる} & \af{＋ 押し売りがない} & \af{− 資料だけでは比べにくい} \\
  接点\newline\scriptsize 人・物・情報 & \af{店のサイト、口コミ} & \af{受付、サインカード} & \af{展示車、価格パネル} & \af{スタッフ、タブレット} & \af{カフェ、遊具} & \af{資料、スタッフ} & \af{見積もりのデータ、メール} \\
\end{wstabh}

\work[60mm]{ワーク2}{AIにペルソナとして店内を歩かせる}
\begin{promptbox}[ペルソナとして店内を歩く]
あなたは次のペルソナです。以下のレイアウトの店舗に、（来店目的）のために来店しました。入店から退店までの行動と、その時々の気持ちを時系列で書いてください。迷う、待たされる、話しかけられて困る、逆に放置されるなど、不満や迷いが生じる場面を必ず含め、その原因も書いてください。ペルソナ：（要約）　レイアウト：（ゾーンの配置と動線を文章で）
\end{promptbox}
\begin{wstabh}{colspec={X[2,l,m] X[3,l,m] X[3,l,m] Q[c,wd=13mm,m]},row{2-Z}={ht=12mm}}
  場面 & 何が起きたか（つまずき） & 原因 & 優先度\newline 高中低 \\
  \af{展示を見る} & \af{展示車に乗ると、キッズスペースの子どもが見えない} & \af{キッズスペースが展示の奥・左側にあり、車内からの視線が届かない} & \ans{高} \\
  \af{入店} & \af{サインカードの意味がわからず、受付で戸惑う} & \af{カードの説明がなく、初めての客には使い方が伝わらない} & \ans{高} \\
  \af{商談} & \af{「相談したい」カードを出したのに、しばらく放置された} & \af{スタッフの待機位置から、カードが見えにくい} & \ans{中} \\
  \af{待ち時間} & \af{見積もりの待ち時間に子どもが飽きる} & \af{待ち時間の目安がわからない} & \ans{中} \\
  \af{退店後} & \af{家で比べるときに、店で見た内容を思い出せない} & \af{資料が紙だけで、試乗の感想などを残せない} & \ans{低} \\
\end{wstabh}

\work[70mm]{ワーク3}{つまずきを解消する改善案}
\begin{promptbox}[改善案を比較する]
次の「つまずき」を解消するレイアウトの改善案を3つ出し、それぞれの利点、欠点、必要なコストの大小、他のペルソナへの影響を表で比較してください。つまずき：（場面と原因）
\end{promptbox}
\begin{wstabh}{colspec={X[2,l,m] X[3,l,m] X[3,l,m] X[2,l,m]},row{2-Z}={ht=15mm}}
  つまずき & レイアウトの改善 & 運用の改善 & 他ペルソナへの影響・コスト \\
  \af{車内から子どもが見えない} & \af{キッズスペースを展示ゾーンの中央に移し、展示車で囲む} & \af{見守り担当を1名置く} & \af{Bには子どもの声が気になる可能性。家具の移動のみで低コスト} \\
  \af{サインカードの意味がわからない} & \af{入口に使い方を示す大きなサインを置く} & \af{受付で一言だけ説明する（「呼ぶまで声はかけません」）} & \af{Bにも有効。コスト小} \\
  \af{カードを出しても放置される} & \af{スタッフの待機カウンターを展示ゾーンの見渡せる位置に} & \af{カードの色で呼び出しを知らせるタブレット連携} & \af{連携は中コスト。まずは配置の変更から} \\
\end{wstabh}

\work[70mm]{比較}{改善前と改善後のレイアウト}
\noindent
\begin{minipage}[t]{0.49\linewidth}\agridbox[改善前]{65mm}{%
  \draw[ans,thick] (2,2) rectangle (77,58);
  \draw[ans,thick,fill=ans!6] (4,4) rectangle (26,30) node[midway]{キッズ};
  \draw[ans,thick,fill=ans!6] (4,34) rectangle (26,56) node[midway]{カフェ};
  \draw[ans,thick,fill=ans!6] (30,20) rectangle (60,44) node[midway]{展示};
  \draw[ans,thick,fill=ans!6] (63,20) rectangle (75,56) node[midway]{商\\談};
  \draw[ans,thick,fill=ans!6] (32,4) rectangle (52,14) node[midway]{受付};
  \draw[ans,thick,dashed] (45,32) -- (15,17) node[pos=0.45,above,sloped]{× 見えない};
}\end{minipage}\hfill
\begin{minipage}[t]{0.49\linewidth}\agridbox[改善後]{65mm}{%
  \draw[ans,thick] (2,2) rectangle (77,58);
  \draw[ans,thick,fill=ans!6] (28,22) rectangle (50,40) node[midway]{キッズ};
  \draw[ans,thick,fill=ans!6] (8,26) rectangle (22,36) node[midway]{車};
  \draw[ans,thick,fill=ans!6] (56,26) rectangle (70,36) node[midway]{車};
  \draw[ans,thick,fill=ans!6] (32,45) rectangle (46,55) node[midway]{車};
  \draw[ans,thick,fill=ans!6] (4,44) rectangle (24,56) node[midway]{カフェ};
  \draw[ans,thick,fill=ans!6] (54,44) rectangle (75,56) node[midway]{商談};
  \draw[ans,thick,fill=ans!6] (30,4) rectangle (50,14) node[midway]{受付・使い方\\サイン};
  \draw[ans,thick,fill=ans!6] (56,8) rectangle (72,16) node[midway]{待機};
}\end{minipage}\par

\begin{nexttime}
  \item 改善後のレイアウト案と、改善の根拠（どのつまずきをどう解消したか）をまとめる
  \item 企業担当者に確認したい点のリスト（現状の来店客の動き、スタッフの運用など）を作る
\end{nexttime}
\abox[企業担当者に確認したい点]{30mm}{・現在、子連れの来店客はどのくらいいて、どこで困っているように見えるか\newline
・見守りや受付の説明に、スタッフを割けるか（平日・休日の人数）\newline
・キッズスペースを展示ゾーンの中央に置くことに、安全上の基準はあるか\newline
・サインカードのような「声をかけない接客」を試したことはあるか}
\aimemoA{ペルソナとしての店内行動の予測と、改善案の比較}
{「あなたはこのペルソナです。入店から退店までの行動と気持ちを時系列で。不満や迷いの場面を必ず含めて」「改善案を3つ、利点・欠点・コスト・他のペルソナへの影響を表で」}
{予測から5つのつまずきを抽出し、優先度は自分たちで付けた。改善案のうち、大がかりな改装案は予算の制約から採用しなかった。}
{AIの予測は仮説なので、子育て中の家族にサインカードの案を見せて感想を聞いた。企業担当者への確認事項にも加えた}

\begin{point}
  \item ジャーニーマップは、行動・感情・接点の3行が埋まり、感情の「−」の場面が具体的に書けているかを見る。
  \item つまずきの「原因」がレイアウトや運用の言葉で書けていると、改善案につながる。「なんとなく不便」は原因になっていない。
  \item AIの予測は仮説にすぎない。実際の観察や企業の知見で確かめる手立て（確認したい点のリスト）が書けているかを評価する。
  \item 改善の前後の図で、どのつまずきが解消されたかを指させるかを確認する。
\end{point}
